\documentclass[11pt]{article}
\usepackage{ochamath}
\subjectcolor{2F5DA8}
\setsheet{線形代数}{連立方程式と掃き出し法　演習}{https://math.ochanote.com/linear-algebra/linear-systems/}
\namelinetrue
\begin{document}
\makesheethead{問題}

\problem[一意解]
$2x+y=7,\ x-y=2$ を掃き出し法で解き、元の式に代入して確認せよ。
\workspace{38mm}

\problem[解の存在]
$x+2y=1,\ 2x+4y=3$ の解の有無を判定し、係数行列と拡大行列の階数を求めよ。
\workspace{38mm}
\newpage

\problem[自由変数]
$x+y+z=2,\ x-y=0$ の一般解を1つのパラメータで表せ。
\workspace{38mm}

\problem[節点電圧]
各節点と接地の間、および節点間のコンダクタンスを1 Sとする。節点1に2 Aを注入したとき、$2V_1-V_2=2,\ -V_1+2V_2=0$ を解け。
\workspace{38mm}
\end{document}
